\section{Design}
\label{sec:design}

\subsection{Overview}

\system is an inline consumer-side verifier, rather than a new policy or a
remote robot attestation service.  It places one trusted authority path beside
the untrusted proposal path and mediates their only route to the actuator
interface.  Figure~\ref{fig:overview} separates three system stages.  During
\emph{construction and qualification}, the monitor compiles $T$ into
$\mathcal G_T$, registers checker/effect vocabularies, and qualifies their
implementations; it does not generate a reference trajectory.  During
\emph{online authorization}, it freezes $O_t^T$ and $Z_t$, receives one source
\actionblock, assesses it, compiles a contract, and conditionally screens a
guard.  During \emph{execution closure}, a fresh authorization gates the sole
dispatch boundary and ordered receipts/effects determine whether the phase may
advance.  Policy-facing bytes remain evidence of what the VLA consumed, not
authority, and \system neither resamples the policy nor relabels $Z_t$ after
seeing the action.

Across these stages, the three obligations in
$\mathsf{ProofAligned}$ remain explicit.  Eligibility requires a bound
checker-relative assessment of the exact source block.  Transaction alignment
requires fresh authorization, ordered exact dispatch, and matching receipts and
effects before phase advance.  Triggered containment requires a guard that
passes the registered identity, restoration, margin, and force conditions.
Later stages preserve earlier identities rather than replacing their verdicts.

The running example makes the layers non-interchangeable.  While the policy
sees ``move to the farthest fixture,'' the trusted branch remains anchored to
soda-to-plate and L1 assesses the exact returned block against that legal
frontier.  L1 rejects a covered hard failure, but a task-progress mismatch can
remain advisory.  If the block proceeds, L2a binds that same canonical block
to one-use dispatch and receipt/effect evidence; a triggered L2b screen covers
joint-side risk.  Thus, L2 can faithfully dispatch and contain a task-divergent
block, while L1 alone establishes neither sink identity nor containment.  The
composition protects three different edges---trusted context to proposal,
proposal to dispatch evidence, and near-boundary dispatch to covered physical
outcome.

\begin{figure*}[t]
\centering
\resizebox{\textwidth}{!}{%
\begin{tikzpicture}[
  font=\scriptsize,
  flow/.style={-{Stealth[length=1.8mm,width=1.35mm]},line width=.62pt,
               draw=black!78,line cap=round,shorten <=1pt,shorten >=1pt},
  attackflow/.style={-{Stealth[length=1.9mm,width=1.45mm]},line width=.72pt,
                     draw=paOrange!94!black,line cap=round,
                     dash pattern=on 2.4pt off 1.6pt,
                     shorten <=1pt,shorten >=1pt},
  trustflow/.style={-{Stealth[length=1.8mm,width=1.35mm]},line width=.68pt,
                    draw=paBlue!88!black,line cap=round,
                    shorten <=1pt,shorten >=1pt},
  evidenceflow/.style={-{Stealth[length=1.7mm,width=1.25mm]},line width=.6pt,
                       draw=paGray,dash pattern=on 2.2pt off 1.5pt,
                       line cap=round,shorten <=1pt,shorten >=1pt},
  untrusted/.style={draw=paOrange!90!black,rounded corners=3pt,
                    fill=paOrange!5,align=center,minimum height=10mm,
                    text width=23mm,inner sep=2pt,dashed},
  source/.style={draw=black,rounded corners=3pt,fill=white,align=center,
                 minimum height=10mm,text width=24mm,inner sep=2pt,
                 line width=.8pt},
  trustedchip/.style={draw=paBlue!85!black,rounded corners=3pt,fill=paBlue!7,
                      align=center,minimum height=7mm,text width=19mm,
                      inner sep=1.5pt},
  selector/.style={draw=paBlue!85!black,rounded corners=3pt,fill=paBlue!7,
                   align=center,minimum height=12mm,text width=23mm,inner sep=2pt},
  c1/.style={draw=paPurple!85!black,rounded corners=3pt,fill=paPurple!8,
             align=center,minimum height=13mm,text width=25mm,inner sep=2pt},
  c2/.style={draw=paGray,rounded corners=3pt,fill=black!3,
             align=center,minimum height=13mm,text width=24mm,inner sep=2pt},
  c3/.style={draw=paGreen!75!black,rounded corners=3pt,fill=paGreen!8,
             align=center,minimum height=13mm,text width=24mm,inner sep=2pt},
  dispatch/.style={draw=paBlue!85!black,rounded corners=3pt,fill=paBlue!6,
                   align=center,minimum height=13mm,text width=26mm,inner sep=2pt},
  external/.style={draw=paGray,rounded corners=3pt,fill=white,align=center,
                   minimum height=13mm,text width=20mm,inner sep=2pt},
  attackpill/.style={draw=paOrange!94!black,rounded corners=4pt,
                     fill=paOrange!7,align=center,minimum height=7.5mm,
                     font=\scriptsize\bfseries,inner sep=1.7pt},
  badge/.style={circle,draw=paGray,fill=white,minimum size=5.0mm,
                inner sep=0pt,line width=.45pt},
  edgelabel/.style={font=\tiny,fill=white,rounded corners=1pt,
                    inner sep=1pt,text=black!78},
  attacklabel/.style={font=\tiny,fill=white,rounded corners=1pt,
                      inner sep=1pt,text=paOrange!94!black},
  fastpill/.style={font=\tiny,draw=black!35,fill=white,
                   rounded corners=2pt,inner xsep=2pt,inner ysep=1pt,
                   text=black!78}
]

% ---------------------------------------------------------------------------
% Threat-model plane: the policy branch remains outside the software TCB.
% ---------------------------------------------------------------------------
\node[attackpill,text width=28mm] (contextattack) at (.8,3.62)
  {Context attacker};
\node[badge,draw=paOrange!94!black] (contextattackicon)
  at ([xshift=1.5mm,yshift=.7mm]contextattack.north west) {};
\pic[scale=.55] at (contextattackicon) {paAttack};

\node[untrusted] (policyview) at (.8,2.15)
  {Policy inputs\\prompt $\bullet$ vision and state $\bullet$ history};
\node[untrusted] (vla) at (3.35,2.15)
  {Frozen \vla\\one proposal ($K=1$)};
\node[source] (action) at (5.95,2.15)
  {\textbf{Source \actionblock}\\canonical $A_t$};

\node[badge,draw=paOrange!90!black] (policyicon)
  at ([xshift=1.7mm,yshift=.8mm]policyview.north west) {};
\pic[scale=.60] at (policyicon) {paEye};
\node[badge,draw=paOrange!90!black] (vlaicon)
  at ([xshift=1.7mm,yshift=.8mm]vla.north west) {};
\pic[scale=.60] at (vlaicon) {paModel};
\node[badge,draw=black] (actionicon)
  at ([xshift=1.7mm,yshift=.8mm]action.north west) {};
\pic[scale=.60] at (actionicon) {paAction};

\draw[attackflow] (contextattack.south)--(policyview.north);
\draw[flow] (policyview)--(vla);
\draw[flow] (vla)--(action);

% The execution attacker from Fig. 1 names the two protected targets.
\node[attackpill,text width=29mm] (stackattack) at (10.75,3.62)
  {Execution attacker};
\node[badge,draw=paOrange!94!black] (stackattackicon)
  at ([xshift=1.5mm,yshift=.7mm]stackattack.north west) {};
\pic[scale=.55] at (stackattackicon) {paAttack};

% ---------------------------------------------------------------------------
% Trusted online path: authority, assessment, transaction, containment, gate.
% ---------------------------------------------------------------------------
\node[trustedchip] (task) at (.55,.32) {Task $T$};
\node[trustedchip] (obs) at (.55,-.38) {Obs. tap $O_t^T$};
\node[font=\tiny\bfseries,text=paBlue!88!black] at (.55,.93)
  {TRUSTED INPUTS};
\node[badge,draw=paBlue!85!black] (taskicon) at (task.west) {};
\pic[scale=.52] at (taskicon) {paTask};
\node[badge,draw=paBlue!85!black] (obsicon) at (obs.west) {};
\pic[scale=.52] at (obsicon) {paEye};

\coordinate (trustedmerge) at (1.78,0);
\node[selector] (select) at (3.35,0)
  {Task graph $\mathcal G_T$\\frozen selector $Z_t$};
\node[c1] (l1) at (5.95,0)
  {\textbf{C1 \lone assess}\\bind $A_t$ to $(Z_t,O_t^T)$\\verdict $S_t$\\reject $\mid$ advise $\mid$ allow};
\node[c2] (contract) at (8.65,0)
  {\textbf{C2 \ltwoa bind}\\$C_t$, epoch, nonce\\one-use $Auth_t$};
\node[c3] (l2b) at (11.30,0)
  {\textbf{C3 \ltwob screen}\\joint trigger gives $\leq2$ guards\\qualify $g_t$\\restore; margin; force};
\node[dispatch] (dispatchnode) at (14.10,0)
  {\textbf{C2 dispatch boundary}\\fresh auth and exact $A_t$\\ordered prefix $\rightarrow$ command};
\node[external] (plant) at (17.05,0)
  {Robot and\\environment};

\node[badge,draw=paBlue!85!black] (selecticon)
  at ([xshift=1mm,yshift=1mm]select.north west) {};
\pic[scale=.52] at (selecticon) {paSelect};
\node[badge,draw=paPurple!85!black] (l1icon)
  at ([xshift=1mm,yshift=1mm]l1.north west) {};
\pic[scale=.52] at (l1icon) {paShield};
\node[badge] (contracticon)
  at ([xshift=1mm,yshift=1mm]contract.north west) {};
\pic[scale=.52] at (contracticon) {paLedger};
\node[badge,draw=paGreen!75!black] (guardicon)
  at ([xshift=1mm,yshift=1mm]l2b.north west) {};
\pic[scale=.52] at (guardicon) {paShield};
\node[badge,draw=paBlue!85!black] (dispatchicon)
  at ([xshift=1mm,yshift=1mm]dispatchnode.north west) {};
\pic[scale=.52] at (dispatchicon) {paDispatch};
\node[badge] (planticon)
  at ([xshift=1mm,yshift=1mm]plant.north west) {};
\pic[scale=.52] at (planticon) {paRobot};

\draw[trustflow] (task.east)--(trustedmerge);
\draw[trustflow] (obs.east)--(trustedmerge);
\draw[trustflow] (trustedmerge)--(select.west);
\draw[trustflow] (select)--(l1);
\draw[flow] (action.south)--(l1.north);
\draw[flow] (l1)--(contract);
\draw[flow] (contract)--(l2b);
\draw[flow] (l2b)--(dispatchnode);
\draw[flow] (dispatchnode)--(plant);

% No-trigger blocks still require the same authorization and dispatch gate.
% A dedicated pill keeps the condition clear of the line.
\node[fastpill] (fastlabel) at (11.55,1.08)
  {no trigger $\rightarrow$ skip C3};
\draw[draw=black!78,line width=.62pt,line cap=round]
  (contract.north)--(contract.north|-fastlabel.west)--(fastlabel.west);
\draw[flow] (fastlabel.east)--(dispatchnode.north|-fastlabel.east)
  --(dispatchnode.north);

% Attack arrows terminate at the exact objects protected by C1 and C2.
\draw[attackflow] (stackattack.south west)
  to[out=-125,in=75] (action.north);
\draw[attackflow] (stackattack.south east)
  to[out=-55,in=85] (dispatchnode.north);
\node[attacklabel] at (7.15,2.76) {tamper or replay};
\node[attacklabel] at (13.02,2.07) {bypass or reorder};

% ---------------------------------------------------------------------------
% Evidence closure: receipts and effects bind back to the same proposal.
% ---------------------------------------------------------------------------
\node[c2,text width=27mm] (observer) at (14.10,-2.15)
  {\textbf{C2 receipt and effect observer}\\ordered $R_t$ and effect window $E_t$};
\node[c2,text width=36mm] (store) at (9.60,-2.15)
  {\textbf{Protected evidence and phase gate}\\proposal $\rightarrow$ auth $\rightarrow$ receipt $\rightarrow$ effect\\advance only after aligned closure};
\node[badge] (observericon)
  at ([xshift=1mm,yshift=1mm]observer.north west) {};
\pic[scale=.52] at (observericon) {paEye};
\node[badge,draw=paBlue!85!black] (storeicon)
  at ([xshift=1mm,yshift=1mm]store.north west) {};
\pic[scale=.52] at (storeicon) {paEvidence};

\draw[evidenceflow] (plant.south)|-(observer.east);
\draw[evidenceflow] (observer)--(store);

% Trust and transaction regions are behind all labels and icons.
\begin{pgfonlayer}{background}
  \node[draw=paBlue!85!black,densely dotted,fill=paBlue!1,
        rounded corners=4pt,
        fit=(task)(obs)(select)(l1)(contract)(l2b)(dispatchnode)
            (observer)(store)(fastlabel),
        inner xsep=3.5mm,inner ysep=3.2mm] (tcb) {};
  \node[draw=paGray,dashed,fill=black!1,rounded corners=4pt,
        fit=(contract)(l2b)(dispatchnode)(observer)(store)(fastlabel),
        inner xsep=2.4mm,inner ysep=2.5mm] (txn) {};
\end{pgfonlayer}

\node[anchor=south west,font=\scriptsize\bfseries,text=paBlue!85!black,
      fill=white,rounded corners=1pt,inner sep=1.2pt]
  at ([xshift=1mm,yshift=1mm]tcb.south west)
  {\system software TCB};
\node[anchor=south east,font=\tiny\bfseries,text=paGray,
      fill=white,rounded corners=1pt,inner sep=1.1pt]
  at ([xshift=-1mm,yshift=1mm]txn.south east)
  {C2 \ltwoa one-proposal transaction};

\end{tikzpicture}

}
\caption{\system separates trusted semantic authority from the potentially
attacked policy branch.  The one-proposal transaction scope spans contract
construction, the conditional L2b screen, fresh authorization,
runtime-checked dispatch, and evidence closure.  The runtime carries the
canonical source-\actionblock
identity across these software artifacts; the guard/controller configuration
is separately audited and is not part of the current Lean identity.}
\label{fig:overview}
\end{figure*}

\subsection{System Components and Evidence Flow}

The authority compiler and selector emit a task graph and state-epoch-bound
$Z_t$ without reading the policy-facing prompt.  The assessor binds $Z_t$, the
matching observation, and exact source block into $S_t$.  The transaction
engine compiles $C_t$, maintains nonces, and conditionally invokes L2b.  The
sole actuator-facing boundary then checks each ordered row, emits receipts
$R_t$, and opens the effect window for $E_t$ before phase advance.

This placement matches the three evidence endpoints: a policy self-report
cannot establish authority, a plausibility check cannot prove which command
the actuator interface accepted, and a later outcome cannot retroactively
authorize dispatch.  \system instead carries one proposal identity forward
while adding evidence at the authority tap, dispatch boundary, and effect
window; it does not compare the online block with a unique reference trajectory.

\subsection{Task-Model Construction}

The selector can emit one of a finite set of task-level skills: \texttt{pick\_up}, \texttt{move}, \texttt{place}, \texttt{release}, \texttt{open}, \texttt{close}, \texttt{actuate}, \texttt{finish}, or \texttt{unknown}.  A task-specific graph restricts which entities and transitions are legal.  For example, a mug-placement task cannot select \texttt{open(drawer)} even if an untrusted model would score it highly.

In our final system, the authority path is
\begin{equation}
\mathsf{BDDL\ goal}\rightarrow\mathsf{task\ graph}
\rightarrow\mathsf{frontier\ selector}\rightarrow Z_t.
\end{equation}
The compiler supports the six goal predicates present in the frozen study:
\texttt{On}, \texttt{In}, \texttt{Open}, \texttt{Close}, \texttt{TurnOn}, and
\texttt{TurnOff}.  Fixed templates expand each predicate into legal subtask
transitions.  For example, \texttt{On(soda,plate)} expands to
\texttt{pick\_up(soda)} $\rightarrow$ \texttt{move(soda,plate)} $\rightarrow$
\texttt{place(soda,plate)} $\rightarrow$ \texttt{release(soda)}.  At each
replanning boundary, a deterministic privileged-geometry finite-state machine
chooses the current frontier using the trusted goal, gripper state, object
state, and auditable distance relations.

This construction avoids treating a VLA semantic score as authorization.  It
also defines the scope precisely: the selector is a benchmark adapter for a
small task language, not a learned planner or camera-only perception system.
An incorrect but syntactically legal BDDL adapter remains a TCB failure.

Each semantic artifact commits to the full context and a canonical selected subtask.  The L1 checker references the artifact digest.  It cannot rename the action's intent after observing $A_t$, and an observation-epoch change requires a new artifact and assessment.

\subsection{C1: L1 Checker-Relative Assessment}

The checker parses translation, rotation, and gripper commands over the full
$10\times7$ chunk.  It does not run a complete dynamics rollout.  Starting
from trusted end-effector position $\hat p_0$, it approximates the translational
prefix by
\begin{equation}
\hat p_i=\hat p_{i-1}+0.05\,\Delta p_i,
\qquad i\in\{1,\ldots,10\}.
\end{equation}
It combines this prefix with rotation/gripper commands, entity geometry,
object/gripper relations, and the frozen task frontier.  The output is a finite
set of motion, task-effect, precondition, and violation atoms plus a
\texttt{known}/\texttt{unknown} status.

Hard risks cover translation or rotation velocity limits, workspace exit,
unexpected-contact neighborhoods, stale state, malformed commands, and
unsupported unknown evidence.  When any hard risk is present, no dispatch
authorization is created.  When no hard risk is present, L1 returns the
unmodified source action.  The result means ``allowed relative to these rules
and this state,'' not ``proved to implement the natural-language task.''  A
downstream command rewrite must invalidate the assessment and obtain a new
contract and authorization.

Task-progress signals are treated differently.  For example, a block that does
not reduce distance to the current target may still be a legitimate detour.
Closing outside a target neighborhood, missing release/close progress, an
expected task-effect miss, or unavailable articulation state therefore
produces an advisory and mandatory next-block replan.  This asymmetry exposes
proxy uncertainty instead of laundering it into a security theorem.  It is
also a limitation: L1 closes the authorization gap only for its hard-risk
vocabulary and provenance checks.

\begin{table}[t]
\centering
\caption{L1 verdict semantics.  Advisory evidence cannot authorize a later
block without a fresh observation and assessment.}
\label{tab:l1-verdicts}
\footnotesize
\begin{tabularx}{\columnwidth}{l>{\raggedright\arraybackslash}X>{\raggedright\arraybackslash}X}
\toprule
Class & Evidence & Runtime semantics \\
\midrule
Hard & Velocity/workspace bound, unexpected contact, stale/malformed object, unsupported unknown & Reject current block; reobserve \\
Advisory & Task progress, release/close progress, expected-effect miss, unavailable articulation state & May dispatch current block; must replan next \\
Allow & Bound evidence with no hard-risk atom & Dispatch exact $A_t$; normal replan \\
\bottomrule
\end{tabularx}
\end{table}

\subsection{C2: L2a One-Use Execution Transaction}

After L1, or directly in the L2-only arm, the consumer first compiles $C_t$.  It creates a fresh authorization only after a no-risk fast path or a triggered L2b screen qualifies.  L2b selects an execution configuration without changing the source-action bytes, consuming an authorization, or retroactively justifying one.  The dispatch boundary verifies the contract, nonce, proposal index, state epoch, action length, and canonical command identity before applying each step.  Receipts bind the authorization and the applied step.  The effect observer opens only after dispatch and closes at the registered observation boundary.  With L1 disabled, the L2-only schema binds the raw policy proposal, state, authorization, and receipts without claiming semantic alignment; it isolates execution identity and containment rather than trusted-task authorization.

Unknown or missing integrity evidence blocks alignment.  Required task-progress evidence may request replanning rather than retroactively altering the command, whereas command/receipt substitution, stale/replayed state, workspace/wrong-contact evidence, and unrecognized unknown effects fail closed.  For arms with L2 enabled, task phase advances only when the transaction is aligned and task completion is observed.

\begin{figure}[t]
\centering
\resizebox{\columnwidth}{!}{\begin{tikzpicture}[
  font=\scriptsize,
  node distance=4.5mm,
  state/.style={draw,rounded corners=1.5pt,align=center,minimum width=31mm,
                minimum height=7.5mm,fill=gray!8,inner sep=2pt},
  success/.style={state,fill=green!9},
  fail/.style={state,fill=red!7},
  flow/.style={-{Latex[length=1.6mm]},semithick},
  reject/.style={-{Latex[length=1.6mm]},semithick,dashed}
]
\node[state] (proposal) {Frozen proposal\\$(Z_t,O_t^T,A_t,\mathit{epoch})$};
\node[state,below=of proposal] (assessed) {Contract compiled\\canonical action digest};
\node[state,below=of assessed] (authorized) {Authorized\\fresh, one-use nonce};
\node[state,below=of authorized] (dispatched) {Dispatched\\runtime-bound step receipt};
\node[success,below=of dispatched] (closed) {Aligned completion\\receipts + effects close};

\node[fail,right=11mm of assessed] (rejected) {Reject / no dispatch\\stale, malformed, substituted};
\node[fail,right=11mm of dispatched] (pending) {Pending / no phase advance\\replay, prefix, effect mismatch};

\draw[flow] (proposal)--node[right]{assess + bind}(assessed);
\draw[flow] (assessed)--node[right]{screen if triggered}(authorized);
\draw[flow] (authorized)--node[right]{consume once}(dispatched);
\draw[flow] (dispatched)--node[right]{observe window}(closed);
\draw[reject] (proposal.east)-|node[pos=.25,above]{bad binding}(rejected.north);
\draw[reject] (assessed.east)--node[above]{invalid / hard risk}(rejected.west);
\draw[reject] (authorized.east)-|node[pos=.28,above]{replay/substitute}(pending.north);
\draw[reject] (dispatched.east)--(pending.west);

\node[below=3.5mm of closed,align=center,text width=39mm] {
  Task phase advances only from the aligned-completion state.};
\end{tikzpicture}
}
\caption{L2a transaction states.  A triggered L2b screen must qualify between
contract construction and authorization; negative evidence cannot advance the
phase, and a consumed authorization cannot be replayed.}
\label{fig:l2a-transaction}
\end{figure}

\subsection{C3: L2b Bounded Guard Screening}

On every step, the runtime computes the minimum across fourteen joint-side margins.  If the minimum is greater than $0.30$\,rad, the source action takes the fast path: no shadow simulator is invoked.  Otherwise, the runtime identifies the dangerous joint sides and constructs at most two temporary virtual-stop configurations.

Every candidate starts from the same simulator snapshot and executes the same source action for a one-step shadow rollout.  Pre-dispatch eligibility requires exact snapshot restoration, fourteen predicted margins of at least $0.15$\,rad, source-action/guard/restore identity, and a MuJoCo generalized-constraint-force proxy no greater than 10,000.  The runtime selects the weakest eligible uniform guard.  Receipts, realized margins, and prediction error are recorded after dispatch and cannot retroactively qualify a candidate.  If none is eligible, the monitor records a fail-closed task failure without dispatch.

The number ``two'' counts guard configurations, not policy candidates.  This keeps $K=1$ for the VLA and prevents a safety mechanism from being misreported as best-of-$K$ policy improvement.

\subsection{End-to-End Decision Procedure}

For every policy call, the monitor executes the following order.  It first
captures $O_t^T$, advances the trusted FSM, and freezes $Z_t$ before invoking
the policy.  It then receives exactly one $A_t$, constructs the L1 assessment,
and stops if a hard verdict is present.  Otherwise it compiles $C_t$ and, near
a registered joint-side trigger, screens at most two guard configurations from
the same restored snapshot.  A qualifying execution receives a fresh one-use
authorization bound to the episode nonce and current proposal;
each low-level step must match the authorized command and state epoch before
dispatch.  Finally, receipts and effects close the transaction.  A task phase
can advance only after aligned completion; an advisory can only influence the
next policy call.  This order prevents post-action relabeling of $Z_t$, replay
of an old authorization, and use of future outcome as present-time evidence.

\subsection{Failure Semantics}

\system distinguishes four outcomes: \emph{hard reject} for covered physical or integrity failure; \emph{advisory replan} for task-progress uncertainty; \emph{fail-closed no-dispatch} when no safe guard qualifies; and \emph{aligned completion} when the transaction and task evidence both close.  The distinctions are retained in the trace rather than collapsed into task success, allowing the implementation to distinguish safe failure, unsafe success, and utility-preserving executions with severe risk steps.
